\begin{abstract}
Cylinder chess joins the a-file to the h-file, so the board has no edges or corners. We study its
king-and-pawn endgames with combinatorial game theory, following Elkies: a passed pawn wins, and the kings are
frozen in a mutual zugzwang worth zero. We repair this model with a passed-pawn test that respects en
passant, and give conditions under which the frozen kings are exact. Under the new
test, an exhaustive census finds no nimber above star with at most six pawns on ordinary boards up to four
files wide, but star two appears with seven, in one position and its color-reversed twin, confirmed by three
programs; on the three-file cylinder six pawns suffice. We also find king-and-pawn positions worth
one-eighth, up-second star and Tiny 1, and test them against real chess: with frozen kings and queen
promotions only, the one-eighth hinges on one side of one end position. Four-man tablebases on the cylinder
and the torus confirm that king and rook can no longer force mate; on the cylinder, opposite-colored bishops
still can. We correct earlier work on cylinder Dawson's
chess: the full ring is a first-player win exactly when the strip three files shorter has Grundy value
zero.
\end{abstract}

\begin{IEEEkeywords}
Chess, combinatorial game theory, endgame tablebases, game solving, nimbers.
\end{IEEEkeywords}
